\scanpage{10}

\uline{Corollaire} Soit $u:H\to G$ un hom de présch en groupes \inserted{sur $S$}, avec $G$ loc. noeth. \inserted{et \deleted{séparé sur $S$}}, $H$ de type fini sur $S$, \deleted{Alors} $H$ \inserted{plat et} satisfaisant la condition R (p. ex. à fibres connexes...). \deleted{\inserted{Si $G$ séparé sur $S$}} Alors $\underline{\mathrm{Cent}}_G(u)$ est représentable par un sous-préschéma \deleted{fermé} de $G$, \deleted{\unsure{Si de|Celui-ci}} fermé si $G$ est séparé sur $S$.

Si $K$ est un sous-groupe de $G$, \deleted{alors}\deleted{\unsure{de t.f.}} \inserted{de t.f. sur \unsure{$K$|$S$}} $\underline{\mathrm{Transp}}_G(\unsure{u|u'},K)$ est représentable, et si $K$ est un sous-préschéma de $G$ (resp. sous-préschéma fermé de $G$) il en est de $=$ de $\underline{\mathrm{Transp}}_G(u,K)$. Enfin, si \deleted{\unsure{en tt|$u$ et $H$}} $K$ est également plat \deleted{\unsure{et}} sur $S$ et satisfait la condition R, \deleted{alors $\mathrm{Transp}_G$} et $u$ un \uline{mono}, alors $\deleted{\mathrm{Transp}_G(u,K)}$ $\mathrm{Transp}.\deleted{\unsure{\mathrm{str}}}_G(H,K)$ est représentable par un \deleted{\unsure{sous-pr|sous-pré}} schéma de type fini sur $G$, et c'est un sous-préschéma (resp sous-préschéma fermé) s'il en est de $=$ de $H$ et $K$ dans $G$.

En particulier, $\underline{\mathrm{Norm}}_G(H)$ est représentable, c'est un sous-préschéma (sous-préschéma fermé) \inserted{de $G$,} si $H$ l'est. \deleted{\illegible{2}}
% archive sheet number: 9
